\chapter*{卷前说明：一座帝国怎样维持，又怎样失去连接}
\addcontentsline{toc}{chapter}{卷前说明：一座帝国怎样维持，又怎样失去连接}

本卷从前二〇六年项羽重划关中、刘邦受封汉中写起，到二二〇年献帝禅位、曹丕建魏为止。它不把汉写成一条由“盛”而“衰”的帝王线。两百余年间，长安和雒阳的朝廷反复面对同一类难题：怎样把军功集团变成可征税的居民，怎样让王国、郡县、边郡和新附地区进入同一套文书，怎样在远征、军费、赈济和宫廷继承之间分配有限的资源。

年代是阅读的骨架，却不是本卷的唯一道路。西汉章节从关中、封国和北方边境的相互制约展开；东汉章节从雒阳复兴、边地供给、地方网络与宫廷争斗的重新缠绕展开。最后四个专题把国家能力、财政边疆、社会宗教和史料争议并置，避免把那些在同一事件中交织的条件拆成互不相干的“政治史”和“文化史”。

正文只在真正深描的转折处设置可点击的中文事件名。其余年代节点帮助读者核对先后，却不被伪装成一串独立场景。数字、疆界、军队路线与人物动机若没有足够互证，本卷宁可写明材料所不能及的地方，也不以确定的口吻填满空白。
